% GENERATED FILE. Edit canonical lessons, then run npm run generate:books.
% canonical-source-hash: fnv1a64:ef6063dc50a97895
% canonical-lessons: RU-C108-police-officer, RU-C108-shop-assistant, RU-C108-farmer, RU-C108-baker, RU-C108-driver

\chapter{Police officer, Shop assistant, Farmer, Baker, Driver}
\label{ch:ru-police-officer-shop-assistant-farmer-baker-driver}
\hlchaptermodality{108}

\begin{chapteropening}
\textbf{By the end of this chapter:} \emph{I can say a police officer, a shop assistant, a farmer, a baker and a driver.}

Complete the last lesson of chapter 108: i can say a police officer, a shop assistant, a farmer, a baker and a driver.
\end{chapteropening}

\section[politséyskiy]{\ru{полицейский} (politséyskiy) — a police officer}

\label{lesson:RU-C108-police-officer}



\begin{quote}
Before the new one: say the Russian for a nurse, then the Russian for a cook.
\end{quote}

\subsection*{You'll want to know: \ru{полицейский}}
\textbf{\ru{полицейский}} (\emph{politséyskiy}) — \textquotedblleft{}a police officer\textquotedblright{}.

\begin{grammarlens}[title={What you've built}]
One more word for this chapter.
\end{grammarlens}

\subsection*{Your turn}
\begin{itemize}
\raggedright
  \item \emph{Say it:} \emph{politséyskiy}
  \item \emph{Say it:} \emph{politséyskiy}, once more
  \item \emph{Say it:} say \emph{politséyskiy}
  \item \emph{Recall:} say the Russian for a nurse, then the Russian for a cook, then say \emph{politséyskiy} again
\end{itemize}

\begin{tcolorbox}[breakable,colback=teal!4,colframe=teal!35!black,title={Before you move on}]
What does \ru{полицейский} mean? (\textquotedblleft{}A police officer\textquotedblright{}.) Say it once more.
\end{tcolorbox}
\section[prodavéts]{\ru{продавец} (prodavéts) — a shop assistant}

\label{lesson:RU-C108-shop-assistant}



\begin{quote}
Before the new one: say the Russian for a cook, then the Russian for a police officer.
\end{quote}

\subsection*{You'll want to know: \ru{продавец}}
\textbf{\ru{продавец}} (\emph{prodavéts}) — \textquotedblleft{}a shop assistant\textquotedblright{}.

\begin{grammarlens}[title={What you've built}]
One more word for this chapter.
\end{grammarlens}

\subsection*{Your turn}
\begin{itemize}
\raggedright
  \item \emph{Say it:} \emph{prodavéts}
  \item \emph{Say it:} \emph{prodavéts}, once more
  \item \emph{Say it:} say \emph{prodavéts}
  \item \emph{Recall:} say the Russian for a cook, then the Russian for a police officer, then say \emph{prodavéts} again
\end{itemize}

\begin{tcolorbox}[breakable,colback=teal!4,colframe=teal!35!black,title={Before you move on}]
What does \ru{продавец} mean? (\textquotedblleft{}A shop assistant\textquotedblright{}.) Say it once more.
\end{tcolorbox}
\section[férmer]{\ru{фермер} (férmer) — a farmer}

\label{lesson:RU-C108-farmer}



\begin{quote}
Before the new one: say the Russian for a police officer, then the Russian for a shop assistant.
\end{quote}

\subsection*{You'll want to know: \ru{фермер}}
\textbf{\ru{фермер}} (\emph{férmer}) — \textquotedblleft{}a farmer\textquotedblright{}.

\begin{grammarlens}[title={What you've built}]
One more word for this chapter.
\end{grammarlens}

\subsection*{Your turn}
\begin{itemize}
\raggedright
  \item \emph{Say it:} \emph{férmer}
  \item \emph{Say it:} \emph{férmer}, once more
  \item \emph{Say it:} say \emph{férmer}
  \item \emph{Recall:} say the Russian for a police officer, then the Russian for a shop assistant, then say \emph{férmer} again
\end{itemize}

\begin{tcolorbox}[breakable,colback=teal!4,colframe=teal!35!black,title={Before you move on}]
What does \ru{фермер} mean? (\textquotedblleft{}A farmer\textquotedblright{}.) Say it once more.
\end{tcolorbox}
\section[pékar]{\ru{пекарь} (pékar) — a baker}

\label{lesson:RU-C108-baker}



\begin{quote}
Before the new one: say the Russian for a shop assistant, then the Russian for a farmer.
\end{quote}

\subsection*{You'll want to know: \ru{пекарь}}
\textbf{\ru{пекарь}} (\emph{pékar}) — \textquotedblleft{}a baker\textquotedblright{}.

\begin{grammarlens}[title={What you've built}]
One more word for this chapter.
\end{grammarlens}

\subsection*{Your turn}
\begin{itemize}
\raggedright
  \item \emph{Say it:} \emph{pékar}
  \item \emph{Say it:} \emph{pékar}, once more
  \item \emph{Say it:} say \emph{pékar}
  \item \emph{Recall:} say the Russian for a shop assistant, then the Russian for a farmer, then say \emph{pékar} again
\end{itemize}

\begin{tcolorbox}[breakable,colback=teal!4,colframe=teal!35!black,title={Before you move on}]
What does \ru{пекарь} mean? (\textquotedblleft{}A baker\textquotedblright{}.) Say it once more.
\end{tcolorbox}
\section[vodítel]{\ru{водитель} (vodítel) — a driver}

\label{lesson:RU-C108-driver}



\begin{quote}
Before the new one: say the Russian for a farmer, then the Russian for a baker.
\end{quote}

\subsection*{You'll want to know: \ru{водитель}}
\textbf{\ru{водитель}} (\emph{vodítel}) — \textquotedblleft{}a driver\textquotedblright{}.

\begin{grammarlens}[title={What you've built}]
One more word for this chapter.
\end{grammarlens}

\subsection*{Your turn}
\begin{itemize}
\raggedright
  \item \emph{Say it:} \emph{vodítel}
  \item \emph{Say it:} \emph{vodítel}, once more
  \item \emph{Say it:} say \emph{vodítel}
  \item \emph{Recall:} say the Russian for a farmer, then the Russian for a baker, then say \emph{vodítel} again
\end{itemize}

\begin{tcolorbox}[breakable,colback=teal!4,colframe=teal!35!black,title={Before you move on}]
What does \ru{водитель} mean? (\textquotedblleft{}A driver\textquotedblright{}.) Say it once more.
\end{tcolorbox}
